% ==============================================
% 電気・情報関係学会北陸支部連合大会
% 講演論文投稿用 pLaTeX2e 版テンプレートファイル(2段組)
% Usage: platex jhes2026_platex_utf8_template.tex
% ==============================================
\documentclass[a4paper,11pt,lualatex,ja=standard]{bxjsarticle}

% 数式
\usepackage{amsmath,amsfonts}
\usepackage{bm}
% 画像
\usepackage{graphicx}
\usepackage{float} % 図の位置を固定
% フローチャート
\usepackage{tikz}
\usetikzlibrary{positioning,arrows.meta}
% 評価グラフ
\usepackage{pgfplots}
\pgfplotsset{compat=1.18}
% url
\usepackage{hyperref}
% ==============================================
% page format
%  textwidth  = paperwidth(210mm) -oddsidemargin(13mm) -evensidemargin(13mm)
%  textheight = paperheight(297mm) -topmargin(20mm-header分) -bottommargin(18mm)
% ==============================================
  \setlength{\paperwidth}{210mm}
  \setlength{\textwidth}{\paperwidth}
  \addtolength{\textwidth}{-13mm} % oddside margin
  \addtolength{\textwidth}{-13mm} % evenside margin
  \setlength{\oddsidemargin}{-1in}
  \addtolength{\oddsidemargin}{13mm} %%%% ここで左右を調節
%
  \setlength{\paperheight}{297mm}
  \setlength{\textheight}{\paperheight}
  \addtolength{\textheight}{-28mm} % 文章領域の高さの調整
  \setlength{\topmargin}{-15mm}    % ここで上下位置を調節
%
  \setlength{\headheight}{0in}
  \setlength{\headsep}{0in}
  \setlength{\columnsep}{10mm}

\makeatletter
\renewcommand{\section}{\@startsection{section}{1}{\z@}%
   {0pt}% 見出し前の余白なし
   {0pt}% 見出し後の余白なし
   {\reset@font\large\bfseries}}   %section見出しの文字サイズをlargeに変更
\renewcommand{\subsection}{\@startsection{subsection}{2}{\z@}%
   {1.5\Cvs \@plus.5\Cvs \@minus.2\Cvs}%
   {.5\Cvs \@plus.3\Cvs}%
   {\reset@font\normalsize\bfseries}} %subsectoin見出しの文字サイズをnormalsizeに変更
\makeatother

\def\nendo{2026}

\def\secraddress{〒923-1292 石川県能美市旭台 1-1\\
北陸先端科学技術大学院大学 先端科学技術研究科内}
\def\secrmail{jhes2026@ml.jaist.ac.jp}

\pagestyle{empty}

\begin{document}
\twocolumn[% -- １段組にしたい場合は，ここをコメントアウトする
% ==============================================
% ページヘッダ(変更禁止)
% ==============================================
%
\begin{center}
{\Large {\nendo}年度電気・情報関係学会北陸支部連合大会}\\
\vspace{-2mm}\rule{\textwidth}{1mm}
\end{center}
% ==============================================
% 原稿はここから
% ==============================================
%
% 講演題目 =====================================
\hspace*{12mm} 
\begin{minipage}{160mm}  %210mm-13mm-13mm-33mm-33mm
\begin{center}    %左づめの場合はflushleft
{\Large \bf
% 依存タスクのための\nobreak
\\ \vspace*{4mm}	% 複数行の場合，強制改行して，行間を調整
依存タスクの重要度とストリーム優先度を考慮した

単一GPU向け静的スケジューリング・資源割当て手法

}
\end{center}
\end{minipage}
%

\vspace*{4mm}	% ここで行間を調整
%
% 著者名（所属） =====================================
\hspace*{33.0mm} %講演番号用空欄
\begin{minipage}{118mm}  %210mm-13mm-13mm-33mm-33mm
\begin{center}    %左づめの場合はflushleft
{\large
小林智輝\dag,高野恵輔\dag,井口寧\dag, \ddagger

\dag 北陸先端科学技術大学院大学, \dagger 金沢大学
% 小林智輝,高野恵輔,\dag, \dagger井口寧（北陸先端科学技術大学院大学）,井口寧（金沢大学）
}
\end{center}
\end{minipage}

\vspace*{6mm}	% ここで行間を調整
] % -- １段組にしたい場合は，ここをコメントアウトする


% 本文 =====================================

\noindent
\section{はじめに}

近年，GPU上で依存関係を持つ複数タスクを効率的に並列実行することが重要である．しかし，Streaming Multiprocessor（SM）資源の競合により重要タスクの実行が遅延し，全体の完了時間が増大する場合がある．
本研究の目的は，タスク重要度とstream優先度を考慮した静的スケジューリングおよびGPU資源割当てにより，makespanを短縮することである．
\section{提案手法}

\begin{figure}[H]
\centering
\begin{tikzpicture}[
  box/.style={draw, rounded corners, align=center, minimum width=20mm, minimum height=9mm, font=\scriptsize},
  line/.style={-{Latex[length=2.0mm]}, thick}
]
\node[box] (level) {レベル化};
\node[box, right=7mm of level] (bl) {最長後続経路長};
\node[box, right=7mm of bl] (prio) {優先度評価};
\node[box, below=8mm of prio] (assign) {タスク割当て};
\node[box, left=8mm of assign] (resource) {SM割当て};

\draw[line] (level) -- (bl);
\draw[line] (bl) -- (prio);
\draw[line] (prio) -- (assign);
\draw[line] (assign) -- (resource);
\end{tikzpicture}
\caption{提案手法の概要フロー}
\label{fig:proposal-flow}
\end{figure}

図1に提案手法の処理フローを示す．まず，入力STGの依存関係を解析してタスクをレベル化し，各タスクから出口タスクまでの最長経路長を算出する．次に，この最長後続経路長をタスク重要度として優先度を評価し，優先度の高いタスクから，各streamのSM数と予測完了時刻に基づいて完了時間が最短となるstreamへ割り当てる．最後に，重要タスクを実行するstream 0へ十分なSMを確保し，残りのSMをGreen Context（GC）側のstreamへ分割する．

既存手法は，主にカーネル特性や資源競合に基づいて割当てを決定するため，タスク間の依存関係やクリティカルパス上の重要度を割当てに直接反映しない．これに対し，本手法は，依存関係から求めたタスク重要度に加え，stream優先度，SM数および予測完了時刻を統合して，タスクの実行順序と資源割当てを決定する点が異なる．

\section{評価}

提案手法の効果を確認するため，依存グラフのレベル幅とタスクの処理時間分布をもとに，stream数とSM配分を決定した．このとき，stream 0は通常streamとして優先的にSMを確保し，残りはGCへ割り当てることで，GPUの資源利用率の向上を図る．
実装上では，タスクごとの処理時間と依存関係を用いて，実行時間を見積もりながら割り当てを選択した．図2は，同一の重負荷入力に対する逐次実行を基準とした速度向上を比較したものである．その結果，GC-only と既存手法を上回り，提案手法が最も高い加速を示した．

\begin{figure}[t]
\centering
\begin{tikzpicture}
\begin{axis}[
    ybar,
    width=0.95\columnwidth,
    height=5.2cm,
    xlabel={手法},
    ylabel={Speedup [x]},
    symbolic x coords={GC-only,既存手法,提案手法},
    xtick=data,
    ymin=0,
    enlarge x limits=0.25,
    bar width=9pt,
    nodes near coords,
    every node near coord/.append style={font=\scriptsize},
]
\addplot coordinates {(GC-only,1.395) (既存手法,1.395) (提案手法,1.396)};
\end{axis}
\end{tikzpicture}
\caption{逐次実行を基準とした速度向上比較}
\end{figure}

\section{まとめ}

本提案手法は，依存関係を持つGPUタスクに対して，DFG構築，レベル化，最長後続経路長に基づく優先度評価，およびSM分割を統合した手法である．見積もり上ではわずかな改善が確認されたが，実測では改善幅が小さい．今後は，タスク割当てとSM配分の最適化により，実測性能の向上を目指す．

% section 4 ----
% \section*{参考文献}
\begin{thebibliography}{99}
\bibitem{cuda-guide}
NVIDIA Corporation, ``CUDA C++ Programming Guide,'' CUDA Toolkit Documentation v13.1, Dec. 2025.

\bibitem{kesco}
Zejia Lin, et.al, ``KeSCo: Compiler-based Kernel Scheduling for Multi-task GPU Applications,'' in 2023 IEEE 41st International Conference on Computer Design (ICCD), 2023.
\end{thebibliography}


% ==============================================
% 原稿はここまで
% ==============================================
\end{document}


% メモ
% グラフ